\documentclass{article}
\usepackage[dutch]{babel}
\usepackage[T1]{fontenc}
\usepackage[utf8]{inputenc}
\usepackage{amsmath}

\title{Inleiding tot de getaltheorie}
\author{P.~de Vries}
\date{2026}

\begin{document}
\maketitle

\section{Inleiding}

De getaltheorie is een van de oudste takken van de wiskunde.  In dit artikel
behandelen wij de verdeling van priemgetallen en het verband met de
Riemann-zetafunctie.

De zetafunctie wordt gedefinieerd als
\[
  \zeta(s) = \sum_{n=1}^{\infty} \frac{1}{n^s}, \quad \text{Re}(s) > 1.
\]

% IJ digraph usage
Het IJsselmeer is vernoemd naar de rivier de IJssel.  De stad IJmuiden ligt
aan het Noordzeekanaal.  In de wiskunde schrijven wij bijectie en injectie
zonder de IJ-combinatie.

\section{De priemgetalstelling}

\begin{theorem}
Zij $\pi(x)$ het aantal priemgetallen kleiner dan of gelijk aan~$x$.  Dan geldt
\[
  \pi(x) \sim \frac{x}{\ln x} \quad (x \to \infty).
\]
\end{theorem}

Het bewijs maakt gebruik van de analytische eigenschappen van~$\zeta(s)$ en
het feit dat $\zeta(s) \ne 0$ voor $\text{Re}(s) = 1$.

\section{Conclusie}

De verdeling van priemgetallen blijft een centraal onderwerp in de moderne
getaltheorie en analyse.

\end{document}
